\section{Introduction}
\label{sec:introduction}

Finite positive unitary fibers are rigid.  If a normalized matrix trace sees
one unitary moment equal to one, faithfulness forces the unitary to be the
identity.  An infinite diffuse trace changes that conclusion: a Haar unitary
$u$ satisfies
\[
 \tau(u^r)=0\qquad(r\in\ZZ\setminus\{0\}).                 \tag{1.1}
\]
The tensor-prime loop could therefore carry a trivial visible component and
an arbitrary amount of Fourier-null Haar mass without changing any nonzero
repetition moment.  SD-C14 tests whether this exact escape also supplies the
missing determinant-visible phase.

The measure theorem is stronger than an example.  Let $\mu$ be any finite
positive measure on the unit circle.  We prove
\[
 \int_\TT z^r\,d\mu(z)=1\quad(r\ge1)
 \quad\Longleftrightarrow\quad
 \mu=\delta_1+c\,m_{\rm H},\quad c\ge0.                  \tag{1.2}
\]
Haar mass is the only positive diffuse freedom.  A normalized state has total
mass one, so it forces $c=0$ and recovers finite moment rigidity.  Exactness
with $c>0$ therefore depends on a nonnormalized finite positive trace.

The escape is real but spectrally sterile for the frozen scalar determinant.
On the fiber
\[
 \cA=\CC\oplus L(\ZZ),\qquad
 W=1\oplus u,\qquad
 \Phi_c(a\oplus x)=a+c\tau(x),                             \tag{1.3}
\]
all nonzero power traces equal one.  Hence the trace-log germ is
\[
 \exp\Phi_c\bigl(\log(1-qW)\bigr)=1-q,                    \tag{1.4}
\]
exactly the determinant with no Haar sector.  The ghost is not merely hard to
detect; it is deleted coefficient by coefficient from this holomorphic
observable.

\begin{figure}[t]
\centering
\begin{tikzpicture}[
 node distance=7mm and 9mm,
 box/.style={draw,rounded corners,align=center,inner sep=5pt,font=\small,minimum height=10mm},
 arr/.style={-{Latex[length=2mm]},thick}]
\node[box,fill=green!10] (mu) {$\mu=\delta_1+c m_{\rm H}$\\unique positive diffuse escape};
\node[box,fill=blue!8,below left=of mu] (hol) {positive moments $=1$\\$D_c(q)=1-q$\\ghost invisible};
\node[box,fill=gray!14,below=of mu] (norm) {normalize mass\\$\widetilde\Phi_c(W^r)=1/(1+c)$\\ledger diluted};
\node[box,fill=orange!12,below right=of mu] (other) {FK/chiral/coupling\\radial, phase-erased,\\or mixed-word polluted};
\draw[arr] (mu)--node[left,font=\scriptsize]{trace-log}(hol);
\draw[arr] (mu)--node[right,font=\scriptsize]{state}(norm);
\draw[arr] (mu)--node[right,font=\scriptsize]{make visible}(other);
\node[draw,dashed,fit=(hol)(norm)(other),inner sep=4pt,label={[font=\scriptsize]below:no-visibility trilemma}] {};
\end{tikzpicture}
\caption{The Fourier-null Haar sector is the only positive diffuse escape.
The scalar analytic determinant deletes it, normalization changes the Euler
coefficient, and alternative visibility mechanisms no longer own the same
holomorphic repetition ledger.}
\label{fig:haar-map}
\end{figure}


Other observables do not repair the loss while retaining the same data type.
The Fuglede--Kadison determinant is
\[
 |1-q|\max(1,|q|)^c,                                      \tag{1.5}
\]
so it remains blind in the Euler disk and becomes $c$-sensitive only through
a positive radial factor.  Natural self-adjointization squares to the
identity and collapses the Haar phase to $\{\pm1\}$.  Adding both $u$ and
$u^{-1}$ as recurrent edges creates trace-visible balanced words such as
$uu^{-1}=e$, which are new mixed closed paths rather than pure $p^r$
repetitions.

The paper establishes four conclusions.
\begin{enumerate}
 \item Equation~(1.2) gives the complete positive-measure classification.
 \item SD-C14 realizes every $c\ge0$ but is normalized only at $c=0$.
 \item Its analytic determinant is exactly Euler-ledger preserving and
       exactly Haar-sector blind.
 \item Finite approximants, density perturbations, matched inventories, and
       recurrent coupling expose delay, instability, proves-too-much, and
       mixed-word boundaries.
\end{enumerate}

The accompanying deterministic audit evaluates cyclic orders
$N=2,\ldots,64$ through repetition $128$, Haar weights
$c\in\{0.25,1,3\}$ at twelve complex points, and $64$ signed Fourier-density
perturbations.  It also compares three $128$-atom inventories at the same
three Haar weights.  Every predicted first leak occurs at the frozen mode;
the maximum Fuglede--Kadison quadrature residual is
$4.45\times10^{-16}$, while all nine inventory controls have exactly zero
analytic-determinant difference and phase range.  These computations audit
the formulas but do not supply target-zero evidence.

All statements remain within Symbolic Dynamics.  Operator-algebraic tools
describe the symbolic fiber and its trace; they do not launch a separate
operator-realization program.  No analytic continuation, functional
equation, divisor count, or fixed self-adjoint generator is claimed.
